\documentclass[11pt]{article}
% jugeo-common.tex — shared preamble for all Judgment Geometry papers
% Include via: % jugeo-common.tex — shared preamble for all Judgment Geometry papers
% Include via: % jugeo-common.tex — shared preamble for all Judgment Geometry papers
% Include via: \input{jugeo-common}

% ── Packages ────────────────────────────────────────────────────────────
\usepackage{amsmath,amssymb,amsthm,mathtools}
\usepackage{tikz-cd}
\usepackage{listings}
\usepackage{xcolor}
\usepackage{xspace}
\usepackage{booktabs}
\usepackage{multirow}
\usepackage{enumitem}
\usepackage[expansion=false]{microtype}
\usepackage{hyperref}
\usepackage{cleveref}
% \usepackage{thmtools}  % disabled: not available in this TeX installation
\usepackage{float}
\usepackage{caption}
\usepackage{subcaption}
\usepackage{tabularx}
\usepackage{stmaryrd}       % \llbracket, \rrbracket
\usepackage{bbm}            % \mathbbm{1}

% ── Page geometry ───────────────────────────────────────────────────────
\usepackage[margin=1in]{geometry}

% ── Theorem environments ────────────────────────────────────────────────
\theoremstyle{plain}
\newtheorem{theorem}{Theorem}[section]
\newtheorem{lemma}[theorem]{Lemma}
\newtheorem{proposition}[theorem]{Proposition}
\newtheorem{corollary}[theorem]{Corollary}

\theoremstyle{definition}
\newtheorem{definition}[theorem]{Definition}
\newtheorem{example}[theorem]{Example}
\newtheorem{construction}[theorem]{Construction}

\theoremstyle{remark}
\newtheorem{remark}[theorem]{Remark}
\newtheorem{notation}[theorem]{Notation}

% ── Python listings style ──────────────────────────────────────────────
\definecolor{jugeoBlue}{HTML}{2563EB}
\definecolor{jugeoGreen}{HTML}{059669}
\definecolor{jugeoRed}{HTML}{DC2626}
\definecolor{jugeoPurple}{HTML}{7C3AED}
\definecolor{jugeoGray}{HTML}{6B7280}
\definecolor{jugeoBg}{HTML}{F8FAFC}

\lstdefinestyle{jugeo-python}{
  language=Python,
  basicstyle=\ttfamily\small,
  keywordstyle=\color{jugeoBlue}\bfseries,
  stringstyle=\color{jugeoGreen},
  commentstyle=\color{jugeoGray}\itshape,
  numberstyle=\tiny\color{jugeoGray},
  backgroundcolor=\color{jugeoBg},
  numbers=left,
  numbersep=6pt,
  frame=single,
  framerule=0.4pt,
  rulecolor=\color{jugeoGray!40},
  breaklines=true,
  breakatwhitespace=true,
  showstringspaces=false,
  tabsize=4,
  xleftmargin=1.5em,
  framexleftmargin=1.5em,
  morekeywords={dataclass,frozen,field,Enum,IntEnum,auto,Any,
                Mapping,Sequence,Optional,match,case},
  literate={→}{{$\to$}}1 {←}{{$\leftarrow$}}1
           {≤}{{$\leq$}}1 {≥}{{$\geq$}}1 {≠}{{$\neq$}}1
           {∈}{{$\in$}}1 {∀}{{$\forall$}}1 {∃}{{$\exists$}}1
           {⊕}{{$\oplus$}}1 {⊖}{{$\ominus$}}1,
  escapeinside={(*@}{@*)},
}
\lstset{style=jugeo-python}

\lstdefinestyle{jugeo-output}{
  basicstyle=\ttfamily\small,
  backgroundcolor=\color{jugeoBg},
  frame=single,
  framerule=0.4pt,
  rulecolor=\color{jugeoGray!40},
  breaklines=true,
  showstringspaces=false,
  xleftmargin=1.5em,
  framexleftmargin=0.5em,
  numbers=none,
}

% ── JuGeo-specific macros ──────────────────────────────────────────────

% Judgment 8-tuple
\newcommand{\J}{\mathcal{J}}
\newcommand{\Jud}[8]{(#1,\,#2,\,#3,\,#4,\,#5,\,#6,\,#7,\,#8)}
\newcommand{\JudShort}[1]{\J_{#1}}

% Components
\newcommand{\coord}{\mathsf{c}}            % coordinate
\newcommand{\prop}{\varphi}                 % proposition
\newcommand{\carrier}{\mathcal{A}}          % carrier type
\newcommand{\evid}{\mathcal{E}}             % evidence bundle
\newcommand{\oblig}{\mathcal{O}}            % obligations
\newcommand{\obstr}{\mathcal{B}}            % obstructions
\newcommand{\trust}{\mathcal{T}}            % trust annotation
\newcommand{\prov}{\Pi}                     % provenance

% Trust algebra
\newcommand{\Talg}{\mathfrak{T}}            % trust ordered algebra
\newcommand{\Eadm}{\mathcal{E}_{\mathrm{adm}}}
\newcommand{\tleq}{\preceq}                 % trust ordering
\newcommand{\tjoin}{\oplus}                 % conservative join
\newcommand{\tatten}{\ominus}               % attenuation
\newcommand{\tpromote}{\uparrow_\pi}        % promotion
\newcommand{\tdemote}{\downarrow_\chi}       % demotion

% Trust tiers
\newcommand{\Tcontra}{\ensuremath{\mathsf{CONTRADICTED}}}
\newcommand{\Tunver}{\ensuremath{\mathsf{UNVERIFIED}}}
\newcommand{\Tcopilot}{\ensuremath{\mathsf{COPILOT}}}
\newcommand{\Toracle}{\ensuremath{\mathsf{ORACLE}}}
\newcommand{\Truntime}{\ensuremath{\mathsf{RUNTIME}}}
\newcommand{\Tsolver}{\ensuremath{\mathsf{SOLVER}}}
\newcommand{\Tproof}{\ensuremath{\mathsf{PROOF}}}

% Site and geometry
\newcommand{\Site}{\mathcal{S}}             % semantic site
\newcommand{\Cov}{\mathrm{Cov}}            % covering family
\newcommand{\Ob}{\mathrm{Ob}}              % objects
\newcommand{\Mor}{\mathrm{Mor}}            % morphisms
\newcommand{\res}{\mathrm{res}}            % restriction
\newcommand{\Cech}{\check{C}}              % Čech complex
\newcommand{\Hcoh}[1]{H^{#1}}             % cohomology group

% Descent
\newcommand{\descent}{\mathrm{desc}}
\newcommand{\glue}{\mathrm{glue}}
\newcommand{\overlap}[2]{#1 \cap #2}

% Semantic moves
\newcommand{\VERIFY}{\textsc{Verify}}
\newcommand{\CONSTRUCT}{\textsc{Construct}}
\newcommand{\NORMALIZE}{\textsc{Normalize}}
\newcommand{\GLUE}{\textsc{Glue}}
\newcommand{\RESTRICT}{\textsc{Restrict}}
\newcommand{\TRANSPORT}{\textsc{Transport}}
\newcommand{\REPAIR}{\textsc{Repair}}
\newcommand{\PROMOTE}{\textsc{Promote}}
\newcommand{\CHALLENGE}{\textsc{Challenge}}
\newcommand{\ESCALATE}{\textsc{Escalate}}

% Fragments
\newcommand{\QFLIA}{\texttt{QF\_LIA}}
\newcommand{\QFLRA}{\texttt{QF\_LRA}}
\newcommand{\QFBV}{\texttt{QF\_BV}}
\newcommand{\QFUF}{\texttt{QF\_UF}}

% Misc
\newcommand{\Python}{\textsc{Python}}
\newcommand{\JuGeo}{\textsc{JuGeo}}
\newcommand{\LEAN}{\textsc{Lean}\,4}
\newcommand{\Fstar}{F$^{\star}$}
\newcommand{\Coq}{\textsc{Coq}}
\newcommand{\Dafny}{\textsc{Dafny}}

% Common shortcuts
\newcommand{\ie}{i.e.\@\xspace}
\newcommand{\eg}{e.g.\@\xspace}
\newcommand{\cf}{cf.\@\xspace}
\newcommand{\wrt}{w.r.t.\@\xspace}

% Evaluation shortcuts
\newcommand{\numcases}{300}
\newcommand{\numspec}{100}
\newcommand{\numeq}{100}
\newcommand{\numbug}{100}
\newcommand{\medtime}{6.5\,ms}
\newcommand{\totaltime}{1.949\,s}


% ── Packages ────────────────────────────────────────────────────────────
\usepackage{amsmath,amssymb,amsthm,mathtools}
\usepackage{tikz-cd}
\usepackage{listings}
\usepackage{xcolor}
\usepackage{xspace}
\usepackage{booktabs}
\usepackage{multirow}
\usepackage{enumitem}
\usepackage[expansion=false]{microtype}
\usepackage{hyperref}
\usepackage{cleveref}
% \usepackage{thmtools}  % disabled: not available in this TeX installation
\usepackage{float}
\usepackage{caption}
\usepackage{subcaption}
\usepackage{tabularx}
\usepackage{stmaryrd}       % \llbracket, \rrbracket
\usepackage{bbm}            % \mathbbm{1}

% ── Page geometry ───────────────────────────────────────────────────────
\usepackage[margin=1in]{geometry}

% ── Theorem environments ────────────────────────────────────────────────
\theoremstyle{plain}
\newtheorem{theorem}{Theorem}[section]
\newtheorem{lemma}[theorem]{Lemma}
\newtheorem{proposition}[theorem]{Proposition}
\newtheorem{corollary}[theorem]{Corollary}

\theoremstyle{definition}
\newtheorem{definition}[theorem]{Definition}
\newtheorem{example}[theorem]{Example}
\newtheorem{construction}[theorem]{Construction}

\theoremstyle{remark}
\newtheorem{remark}[theorem]{Remark}
\newtheorem{notation}[theorem]{Notation}

% ── Python listings style ──────────────────────────────────────────────
\definecolor{jugeoBlue}{HTML}{2563EB}
\definecolor{jugeoGreen}{HTML}{059669}
\definecolor{jugeoRed}{HTML}{DC2626}
\definecolor{jugeoPurple}{HTML}{7C3AED}
\definecolor{jugeoGray}{HTML}{6B7280}
\definecolor{jugeoBg}{HTML}{F8FAFC}

\lstdefinestyle{jugeo-python}{
  language=Python,
  basicstyle=\ttfamily\small,
  keywordstyle=\color{jugeoBlue}\bfseries,
  stringstyle=\color{jugeoGreen},
  commentstyle=\color{jugeoGray}\itshape,
  numberstyle=\tiny\color{jugeoGray},
  backgroundcolor=\color{jugeoBg},
  numbers=left,
  numbersep=6pt,
  frame=single,
  framerule=0.4pt,
  rulecolor=\color{jugeoGray!40},
  breaklines=true,
  breakatwhitespace=true,
  showstringspaces=false,
  tabsize=4,
  xleftmargin=1.5em,
  framexleftmargin=1.5em,
  morekeywords={dataclass,frozen,field,Enum,IntEnum,auto,Any,
                Mapping,Sequence,Optional,match,case},
  literate={→}{{$\to$}}1 {←}{{$\leftarrow$}}1
           {≤}{{$\leq$}}1 {≥}{{$\geq$}}1 {≠}{{$\neq$}}1
           {∈}{{$\in$}}1 {∀}{{$\forall$}}1 {∃}{{$\exists$}}1
           {⊕}{{$\oplus$}}1 {⊖}{{$\ominus$}}1,
  escapeinside={(*@}{@*)},
}
\lstset{style=jugeo-python}

\lstdefinestyle{jugeo-output}{
  basicstyle=\ttfamily\small,
  backgroundcolor=\color{jugeoBg},
  frame=single,
  framerule=0.4pt,
  rulecolor=\color{jugeoGray!40},
  breaklines=true,
  showstringspaces=false,
  xleftmargin=1.5em,
  framexleftmargin=0.5em,
  numbers=none,
}

% ── JuGeo-specific macros ──────────────────────────────────────────────

% Judgment 8-tuple
\newcommand{\J}{\mathcal{J}}
\newcommand{\Jud}[8]{(#1,\,#2,\,#3,\,#4,\,#5,\,#6,\,#7,\,#8)}
\newcommand{\JudShort}[1]{\J_{#1}}

% Components
\newcommand{\coord}{\mathsf{c}}            % coordinate
\newcommand{\prop}{\varphi}                 % proposition
\newcommand{\carrier}{\mathcal{A}}          % carrier type
\newcommand{\evid}{\mathcal{E}}             % evidence bundle
\newcommand{\oblig}{\mathcal{O}}            % obligations
\newcommand{\obstr}{\mathcal{B}}            % obstructions
\newcommand{\trust}{\mathcal{T}}            % trust annotation
\newcommand{\prov}{\Pi}                     % provenance

% Trust algebra
\newcommand{\Talg}{\mathfrak{T}}            % trust ordered algebra
\newcommand{\Eadm}{\mathcal{E}_{\mathrm{adm}}}
\newcommand{\tleq}{\preceq}                 % trust ordering
\newcommand{\tjoin}{\oplus}                 % conservative join
\newcommand{\tatten}{\ominus}               % attenuation
\newcommand{\tpromote}{\uparrow_\pi}        % promotion
\newcommand{\tdemote}{\downarrow_\chi}       % demotion

% Trust tiers
\newcommand{\Tcontra}{\ensuremath{\mathsf{CONTRADICTED}}}
\newcommand{\Tunver}{\ensuremath{\mathsf{UNVERIFIED}}}
\newcommand{\Tcopilot}{\ensuremath{\mathsf{COPILOT}}}
\newcommand{\Toracle}{\ensuremath{\mathsf{ORACLE}}}
\newcommand{\Truntime}{\ensuremath{\mathsf{RUNTIME}}}
\newcommand{\Tsolver}{\ensuremath{\mathsf{SOLVER}}}
\newcommand{\Tproof}{\ensuremath{\mathsf{PROOF}}}

% Site and geometry
\newcommand{\Site}{\mathcal{S}}             % semantic site
\newcommand{\Cov}{\mathrm{Cov}}            % covering family
\newcommand{\Ob}{\mathrm{Ob}}              % objects
\newcommand{\Mor}{\mathrm{Mor}}            % morphisms
\newcommand{\res}{\mathrm{res}}            % restriction
\newcommand{\Cech}{\check{C}}              % Čech complex
\newcommand{\Hcoh}[1]{H^{#1}}             % cohomology group

% Descent
\newcommand{\descent}{\mathrm{desc}}
\newcommand{\glue}{\mathrm{glue}}
\newcommand{\overlap}[2]{#1 \cap #2}

% Semantic moves
\newcommand{\VERIFY}{\textsc{Verify}}
\newcommand{\CONSTRUCT}{\textsc{Construct}}
\newcommand{\NORMALIZE}{\textsc{Normalize}}
\newcommand{\GLUE}{\textsc{Glue}}
\newcommand{\RESTRICT}{\textsc{Restrict}}
\newcommand{\TRANSPORT}{\textsc{Transport}}
\newcommand{\REPAIR}{\textsc{Repair}}
\newcommand{\PROMOTE}{\textsc{Promote}}
\newcommand{\CHALLENGE}{\textsc{Challenge}}
\newcommand{\ESCALATE}{\textsc{Escalate}}

% Fragments
\newcommand{\QFLIA}{\texttt{QF\_LIA}}
\newcommand{\QFLRA}{\texttt{QF\_LRA}}
\newcommand{\QFBV}{\texttt{QF\_BV}}
\newcommand{\QFUF}{\texttt{QF\_UF}}

% Misc
\newcommand{\Python}{\textsc{Python}}
\newcommand{\JuGeo}{\textsc{JuGeo}}
\newcommand{\LEAN}{\textsc{Lean}\,4}
\newcommand{\Fstar}{F$^{\star}$}
\newcommand{\Coq}{\textsc{Coq}}
\newcommand{\Dafny}{\textsc{Dafny}}

% Common shortcuts
\newcommand{\ie}{i.e.\@\xspace}
\newcommand{\eg}{e.g.\@\xspace}
\newcommand{\cf}{cf.\@\xspace}
\newcommand{\wrt}{w.r.t.\@\xspace}

% Evaluation shortcuts
\newcommand{\numcases}{300}
\newcommand{\numspec}{100}
\newcommand{\numeq}{100}
\newcommand{\numbug}{100}
\newcommand{\medtime}{6.5\,ms}
\newcommand{\totaltime}{1.949\,s}


% ── Packages ────────────────────────────────────────────────────────────
\usepackage{amsmath,amssymb,amsthm,mathtools}
\usepackage{tikz-cd}
\usepackage{listings}
\usepackage{xcolor}
\usepackage{xspace}
\usepackage{booktabs}
\usepackage{multirow}
\usepackage{enumitem}
\usepackage[expansion=false]{microtype}
\usepackage{hyperref}
\usepackage{cleveref}
% \usepackage{thmtools}  % disabled: not available in this TeX installation
\usepackage{float}
\usepackage{caption}
\usepackage{subcaption}
\usepackage{tabularx}
\usepackage{stmaryrd}       % \llbracket, \rrbracket
\usepackage{bbm}            % \mathbbm{1}

% ── Page geometry ───────────────────────────────────────────────────────
\usepackage[margin=1in]{geometry}

% ── Theorem environments ────────────────────────────────────────────────
\theoremstyle{plain}
\newtheorem{theorem}{Theorem}[section]
\newtheorem{lemma}[theorem]{Lemma}
\newtheorem{proposition}[theorem]{Proposition}
\newtheorem{corollary}[theorem]{Corollary}

\theoremstyle{definition}
\newtheorem{definition}[theorem]{Definition}
\newtheorem{example}[theorem]{Example}
\newtheorem{construction}[theorem]{Construction}

\theoremstyle{remark}
\newtheorem{remark}[theorem]{Remark}
\newtheorem{notation}[theorem]{Notation}

% ── Python listings style ──────────────────────────────────────────────
\definecolor{jugeoBlue}{HTML}{2563EB}
\definecolor{jugeoGreen}{HTML}{059669}
\definecolor{jugeoRed}{HTML}{DC2626}
\definecolor{jugeoPurple}{HTML}{7C3AED}
\definecolor{jugeoGray}{HTML}{6B7280}
\definecolor{jugeoBg}{HTML}{F8FAFC}

\lstdefinestyle{jugeo-python}{
  language=Python,
  basicstyle=\ttfamily\small,
  keywordstyle=\color{jugeoBlue}\bfseries,
  stringstyle=\color{jugeoGreen},
  commentstyle=\color{jugeoGray}\itshape,
  numberstyle=\tiny\color{jugeoGray},
  backgroundcolor=\color{jugeoBg},
  numbers=left,
  numbersep=6pt,
  frame=single,
  framerule=0.4pt,
  rulecolor=\color{jugeoGray!40},
  breaklines=true,
  breakatwhitespace=true,
  showstringspaces=false,
  tabsize=4,
  xleftmargin=1.5em,
  framexleftmargin=1.5em,
  morekeywords={dataclass,frozen,field,Enum,IntEnum,auto,Any,
                Mapping,Sequence,Optional,match,case},
  literate={→}{{$\to$}}1 {←}{{$\leftarrow$}}1
           {≤}{{$\leq$}}1 {≥}{{$\geq$}}1 {≠}{{$\neq$}}1
           {∈}{{$\in$}}1 {∀}{{$\forall$}}1 {∃}{{$\exists$}}1
           {⊕}{{$\oplus$}}1 {⊖}{{$\ominus$}}1,
  escapeinside={(*@}{@*)},
}
\lstset{style=jugeo-python}

\lstdefinestyle{jugeo-output}{
  basicstyle=\ttfamily\small,
  backgroundcolor=\color{jugeoBg},
  frame=single,
  framerule=0.4pt,
  rulecolor=\color{jugeoGray!40},
  breaklines=true,
  showstringspaces=false,
  xleftmargin=1.5em,
  framexleftmargin=0.5em,
  numbers=none,
}

% ── JuGeo-specific macros ──────────────────────────────────────────────

% Judgment 8-tuple
\newcommand{\J}{\mathcal{J}}
\newcommand{\Jud}[8]{(#1,\,#2,\,#3,\,#4,\,#5,\,#6,\,#7,\,#8)}
\newcommand{\JudShort}[1]{\J_{#1}}

% Components
\newcommand{\coord}{\mathsf{c}}            % coordinate
\newcommand{\prop}{\varphi}                 % proposition
\newcommand{\carrier}{\mathcal{A}}          % carrier type
\newcommand{\evid}{\mathcal{E}}             % evidence bundle
\newcommand{\oblig}{\mathcal{O}}            % obligations
\newcommand{\obstr}{\mathcal{B}}            % obstructions
\newcommand{\trust}{\mathcal{T}}            % trust annotation
\newcommand{\prov}{\Pi}                     % provenance

% Trust algebra
\newcommand{\Talg}{\mathfrak{T}}            % trust ordered algebra
\newcommand{\Eadm}{\mathcal{E}_{\mathrm{adm}}}
\newcommand{\tleq}{\preceq}                 % trust ordering
\newcommand{\tjoin}{\oplus}                 % conservative join
\newcommand{\tatten}{\ominus}               % attenuation
\newcommand{\tpromote}{\uparrow_\pi}        % promotion
\newcommand{\tdemote}{\downarrow_\chi}       % demotion

% Trust tiers
\newcommand{\Tcontra}{\ensuremath{\mathsf{CONTRADICTED}}}
\newcommand{\Tunver}{\ensuremath{\mathsf{UNVERIFIED}}}
\newcommand{\Tcopilot}{\ensuremath{\mathsf{COPILOT}}}
\newcommand{\Toracle}{\ensuremath{\mathsf{ORACLE}}}
\newcommand{\Truntime}{\ensuremath{\mathsf{RUNTIME}}}
\newcommand{\Tsolver}{\ensuremath{\mathsf{SOLVER}}}
\newcommand{\Tproof}{\ensuremath{\mathsf{PROOF}}}

% Site and geometry
\newcommand{\Site}{\mathcal{S}}             % semantic site
\newcommand{\Cov}{\mathrm{Cov}}            % covering family
\newcommand{\Ob}{\mathrm{Ob}}              % objects
\newcommand{\Mor}{\mathrm{Mor}}            % morphisms
\newcommand{\res}{\mathrm{res}}            % restriction
\newcommand{\Cech}{\check{C}}              % Čech complex
\newcommand{\Hcoh}[1]{H^{#1}}             % cohomology group

% Descent
\newcommand{\descent}{\mathrm{desc}}
\newcommand{\glue}{\mathrm{glue}}
\newcommand{\overlap}[2]{#1 \cap #2}

% Semantic moves
\newcommand{\VERIFY}{\textsc{Verify}}
\newcommand{\CONSTRUCT}{\textsc{Construct}}
\newcommand{\NORMALIZE}{\textsc{Normalize}}
\newcommand{\GLUE}{\textsc{Glue}}
\newcommand{\RESTRICT}{\textsc{Restrict}}
\newcommand{\TRANSPORT}{\textsc{Transport}}
\newcommand{\REPAIR}{\textsc{Repair}}
\newcommand{\PROMOTE}{\textsc{Promote}}
\newcommand{\CHALLENGE}{\textsc{Challenge}}
\newcommand{\ESCALATE}{\textsc{Escalate}}

% Fragments
\newcommand{\QFLIA}{\texttt{QF\_LIA}}
\newcommand{\QFLRA}{\texttt{QF\_LRA}}
\newcommand{\QFBV}{\texttt{QF\_BV}}
\newcommand{\QFUF}{\texttt{QF\_UF}}

% Misc
\newcommand{\Python}{\textsc{Python}}
\newcommand{\JuGeo}{\textsc{JuGeo}}
\newcommand{\LEAN}{\textsc{Lean}\,4}
\newcommand{\Fstar}{F$^{\star}$}
\newcommand{\Coq}{\textsc{Coq}}
\newcommand{\Dafny}{\textsc{Dafny}}

% Common shortcuts
\newcommand{\ie}{i.e.\@\xspace}
\newcommand{\eg}{e.g.\@\xspace}
\newcommand{\cf}{cf.\@\xspace}
\newcommand{\wrt}{w.r.t.\@\xspace}

% Evaluation shortcuts
\newcommand{\numcases}{300}
\newcommand{\numspec}{100}
\newcommand{\numeq}{100}
\newcommand{\numbug}{100}
\newcommand{\medtime}{6.5\,ms}
\newcommand{\totaltime}{1.949\,s}


\title{\textbf{An Algebraic Foundation for Evidence-Carrying Proofs}}

\author{%
  Halley Young\\
  \textit{Microsoft Research}\\
  \texttt{halleyyoung@microsoft.com}
}

\date{\today}

\begin{document}
\maketitle

\begin{abstract}
We introduce the \emph{judgment}
$\J = (\coord,\,\prop,\,\carrier,\,\evid,\,\oblig,\,\obstr,\,\trust,\,\prov)$,
the central algebraic object of Judgment Geometry: an eight-component
structure that extends the classical proof-theoretic judgment with
evidence bundles, trust annotations, and persistent obstructions.
Traditional type-theoretic judgments
(Martin-L\"of's $\Gamma \vdash a : A$, the Calculus of Inductive Constructions,
\Fstar{} computation types)
record \emph{what} was proved and \emph{where}, but discard information about
\emph{who} proved it, \emph{how much} the proof should be trusted, \emph{what}
obligations remain, and \emph{what} obstructions block completion.
Our formulation extends the classical three-component judgment with evidence bundles,
a trust annotation drawn from an ordered algebra, residual obligations,
persistent obstructions, and a cryptographic provenance trail.
We define five algebraic operations on judgments---restriction, transport,
composition, comparison, and join---and prove the standard structural rules
(weakening, contraction, exchange, cut) admissible.
We implement the algebra in \Python{} and evaluate on 100 benchmark
programs across 10 domains, verifying all 7 trust algebra axioms on every
program (700 axiom checks, 100\% pass rate).  The algebra preserves all
eight fields through every operation---a $2.0\times$ information advantage
over the best competing system (\Fstar{}, which retains four of eight
components).  Each operation completes in $0.15$--$19.5$\,\textmu{}s,
confirming negligible runtime overhead.
\end{abstract}

\tableofcontents
\newpage

%% ═══════════════════════════════════════════════════════════════════════
\section{Introduction}\label{sec:intro}
%% ═══════════════════════════════════════════════════════════════════════

\subsection{Judgments in type theory}

Every proof assistant rests on a notion of \emph{judgment}.
In Martin-L\"of type theory~\cite{martinlof84}, a judgment takes the form
\[
  \Gamma \;\vdash\; a : A,
\]
asserting that term~$a$ inhabits type~$A$ under context~$\Gamma$.
The judgment carries exactly three data: a context, a term, and a type.
The Calculus of Inductive Constructions (CIC), which underlies \Coq{} and
\LEAN{}, enriches the context with universe levels but retains the same
three-component shape~\cite{coquand88,lean4}.

\Fstar{}~\cite{fstar} extends judgments to four components by indexing
computation types with a \emph{weakest-precondition} transformer:
\[
  \Gamma \;\vdash\; e : \mathsf{M}\;a\;\mathit{wp},
\]
where $\mathsf{M}$ is an effect label and $\mathit{wp}$ encodes the
specification.  \Dafny{}~\cite{dafny} threads pre- and postconditions
through a Hoare-style judgment but likewise records no metadata beyond
the logical assertion.

\subsection{What existing judgments discard}

All of these systems answer \emph{what} was proved and \emph{in which context},
but they systematically discard at least five categories of information that
matter in practice:

\begin{enumerate}[label=(\roman*)]
  \item \textbf{Evidence provenance.}  Was the proof found by a human, an SMT
        solver, a runtime witness, or a language-model oracle?  Existing
        systems erase this distinction the moment a proof term is type-checked.
  \item \textbf{Trust gradation.}  A proof discharged by Z3 in 2\,ms and a
        proof reviewed by three experts for a week receive the same
        \texttt{Qed}.
  \item \textbf{Residual obligations.}  Partial proofs---common in
        large-scale verification---have no first-class status; they are
        either admitted (unsound) or rejected (unusable).
  \item \textbf{Obstructions.}  When a proof attempt \emph{fails}, the failure
        is not recorded in the judgment; it vanishes.  In \JuGeo{}, failures
        are first-class cohomology classes.
  \item \textbf{Audit trail.}  No mainstream system maintains a
        machine-readable record of \emph{how} a judgment was derived---which
        solver was invoked, which lemmas were used, which rewrites were
        applied.
\end{enumerate}

\subsection{Contributions}

We propose the \emph{judgment 8-tuple} as a uniform solution.  Our
contributions are:
\begin{enumerate}
  \item A formal definition of $\J = (\coord,\prop,\carrier,\evid,\oblig,
        \obstr,\trust,\prov)$ with well-formedness conditions
        (\S\ref{sec:definition}).
  \item Five algebraic operations---restriction, transport, composition,
        comparison, and join---with functoriality and coherence laws
        (\S\ref{sec:algebra}).
  \item Proofs that the classical structural rules (weakening, contraction,
        exchange, cut) are admissible (\S\ref{sec:structural}).
  \item Logical rules (introduction, elimination, computation) as judgment
        constructors (\S\ref{sec:logical}).
  \item A judgment transition system with small-step semantics
        (\S\ref{sec:transitions}).
  \item An implementation in \Python{} with evaluation on 100 programs verifying
        7 trust algebra axioms (700 axiom checks)
        (\S\ref{sec:evaluation}).
  \item An embedding of CIC, \Fstar{}, and Hoare-triple judgments as
        degenerate 8-tuples (\S\ref{sec:recovering}).
  \item Metatheoretic results: subject reduction, cut elimination, and
        trust monotonicity (\S\ref{sec:metatheory}).
\end{enumerate}


%% ═══════════════════════════════════════════════════════════════════════
\section{The 8-Tuple: Formal Definition}\label{sec:definition}
%% ═══════════════════════════════════════════════════════════════════════

\subsection{Motivation}

The judgment 8-tuple is the central object of Judgment Geometry.
It is implemented in the \JuGeo{} codebase as the \texttt{Judgment} dataclass:

\begin{lstlisting}[caption={The Judgment dataclass (from \texttt{judgment\_terms.py}).},
  label=lst:judgment]
@dataclass(frozen=True, slots=True)
class Judgment:
    """J = (c, (*@$\varphi$@*), A, E, O, B, T, (*@$\Pi$@*)) -- NEVER a boolean."""
    coordinate:   CoordinateObject            # c
    proposition:  Proposition                 # (*@$\varphi$@*)
    carrier:      Carrier                     # A
    evidence:     EvidenceBundle              # E
    obligations:  tuple[ResidualObligation, ...]  # O
    obstructions: tuple[Obstruction, ...]     # B
    trust:        TrustAnnotation             # T
    provenance:   Provenance                  # (*@$\Pi$@*)
\end{lstlisting}

The dataclass is \emph{frozen}: once constructed, a judgment is immutable.
Every operation that ``modifies'' a judgment returns a fresh copy.
This ensures referential transparency and simplifies the metatheory.

\subsection{Component specifications}

We now define each component formally.

\begin{definition}[Judgment 8-tuple]\label{def:8tuple}
A \emph{judgment} is an element
\[
  \J \;=\; \Jud{\coord}{\prop}{\carrier}{\evid}{\oblig}{\obstr}{\trust}{\prov}
\]
where the components are drawn from the following domains:
\begin{enumerate}[label=(\alph*)]
  \item $\coord \in \Ob(\Site)$: a \emph{coordinate} in the semantic site
        $\Site$.  Coordinates index the ``where'' of a judgment---a function,
        a module, a specification point.
  \item $\prop \in \mathrm{Prop}(\Site)$: a \emph{proposition} classified by
        kind $\kappa \in \{\mathsf{structural},\, \mathsf{behavioral},\,
        \mathsf{relational},\, \mathsf{resource},\, \mathsf{semantic}\}$.
  \item $\carrier$: a \emph{carrier type} (the ``$A$'' of $\Gamma \vdash a : A$)
        equipped with a refinement lattice $(\carrier, \sqsubseteq, \sqcap,
        \sqcup)$.
  \item $\evid = (e_1, \ldots, e_n)$: an \emph{evidence bundle}, where each
        $e_i$ carries a trust level, a channel, a timestamp, and an optional
        expiry.
  \item $\oblig = (o_1, \ldots, o_m)$: \emph{residual obligations}---proof
        goals that remain undischarged.
  \item $\obstr = (b_1, \ldots, b_k)$: \emph{obstructions}---persistent
        failure records, modeled as $H^1$ cohomology classes.
  \item $\trust \in \Talg$: a \emph{trust annotation} drawn from the
        ordered algebra $(\Eadm, \tleq, \tjoin, \tatten, \tpromote,
        \tdemote)$.
  \item $\prov$: a \emph{provenance} record containing the derivation source,
        parent judgments, and transformation history.
\end{enumerate}
\end{definition}

\subsubsection{The trust algebra}

The trust component is not a scalar.  It is a structured element of the
\emph{admissible trust algebra} $\Talg = (\Eadm, \tleq, \tjoin, \tatten,
\tpromote, \tdemote)$, where $\Eadm$ is the set of admissible trust
configurations, $\tleq$ is a partial order, $\tjoin$ is a conservative join,
$\tatten$ is an attenuation (demotion) operator, $\tpromote$ is a
policy-bounded promotion, and $\tdemote$ is a challenge-triggered demotion.

The implementation defines six discrete trust levels:

\begin{lstlisting}[caption={Trust levels (from \texttt{judgment\_terms.py}).},
  label=lst:trust]
class TrustLevel(IntEnum):
    CONTRADICTED      = 0   # Evidence contradicts the claim
    UNVERIFIED        = 1   # No evidence yet
    COPILOT_SUGGESTED = 2   # LLM oracle proposed
    RUNTIME_WITNESSED = 3   # Runtime execution confirmed
    SOLVER_DISCHARGED = 4   # SMT solver proved
    VERIFIED_PROOF    = 5   # Formal proof checked
\end{lstlisting}

These levels form a total order
$\Tcontra < \Tunver < \Tcopilot < \Truntime < \Tsolver < \Tproof$.
The conservative join $\tjoin$ takes the \emph{minimum} of two trust levels
(not the maximum), ensuring that a bundle is only as trusted as its weakest
link.

\subsubsection{Judgment status}

A judgment progresses through a lifecycle:

\begin{lstlisting}[caption={Judgment status lifecycle.}, label=lst:status]
class JudgmentStatus(str, Enum):
    PROPOSED    = "proposed"     # Claim made, not yet verified
    CHALLENGED  = "challenged"  # Counter-evidence presented
    SETTLED     = "settled"     # All obligations discharged
    OBSTRUCTED  = "obstructed"  # Persistent obstruction found
\end{lstlisting}

The status is \emph{derived} from the components: a judgment is
$\mathsf{SETTLED}$ iff $\oblig = ()$ and $\obstr = ()$;
$\mathsf{OBSTRUCTED}$ iff $\obstr \neq ()$; and so on.

\subsubsection{Evidence items}

Each evidence item records both its provenance and its epistemic weight:

\begin{lstlisting}[caption={Evidence item classification.}, label=lst:evidence]
class EvidenceItemKind(str, Enum):
    SOLVER_PROOF     = "solver_proof"      # Z3/CVC5 certificate
    RUNTIME_WITNESS  = "runtime_witness"   # Execution trace
    ORACLE_PROPOSAL  = "oracle_proposal"   # LLM suggestion
    FORMAL_PROOF     = "formal_proof"      # Verified proof term
\end{lstlisting}

\subsubsection{Proposition kinds}

Propositions are classified by the \emph{kind} of property they express:

\begin{lstlisting}[caption={Proposition kinds.}, label=lst:propkind]
class PropositionKind(str, Enum):
    STRUCTURAL  = "structural"   # Type-level properties
    BEHAVIORAL  = "behavioral"   # Input-output specifications
    RELATIONAL  = "relational"   # Equivalence / refinement
    RESOURCE    = "resource"     # Memory, time, allocation
    SEMANTIC    = "semantic"     # Domain-specific meaning
\end{lstlisting}

\subsection{Well-formedness conditions}

\begin{definition}[Well-formed judgment]\label{def:wf}
A judgment $\J = \Jud{\coord}{\prop}{\carrier}{\evid}{\oblig}{\obstr}
{\trust}{\prov}$ is \emph{well-formed} if:
\begin{enumerate}[label=(WF\arabic*)]
  \item \label{wf:coord} $\coord \in \Ob(\Site)$ is a valid coordinate.
  \item \label{wf:prop} $\prop$ is a syntactically well-formed proposition
        with $\mathrm{fv}(\prop) \subseteq \mathrm{dom}(\coord)$.
  \item \label{wf:trust-ev} The trust annotation is \emph{consistent} with
        the evidence:
        $\trust.\mathsf{level} \leq \min_{e \in \evid} e.\mathsf{trust\_level}$.
  \item \label{wf:oblig} Every obligation $o_i \in \oblig$ references
        propositions whose free variables are bound by~$\coord$.
  \item \label{wf:obstr} Every obstruction $b_j \in \obstr$ has a non-empty
        description.
  \item \label{wf:prov} The provenance $\prov$ records at least one source.
\end{enumerate}
\end{definition}

\subsection{The judgment as a fiber}

Categorically, a judgment is a \emph{fiber} over its coordinate in the
semantic site.  The functor
\[
  \J(-) : \Site^{\mathrm{op}} \longrightarrow \mathbf{Set}
\]
assigns to each coordinate $U \in \Ob(\Site)$ the set of well-formed
judgments at~$U$.  A morphism $f : V \to U$ in $\Site$ induces a
\emph{restriction} $\J(U) \to \J(V)$, making $\J(-)$ a presheaf.

\begin{equation}\label{eq:fiber}
\begin{tikzcd}[column sep=large, row sep=large]
  \J(U) \arrow[d, "\pi_U"'] \arrow[r, "\res_f"] & \J(V) \arrow[d, "\pi_V"] \\
  U & V \arrow[l, "f"]
\end{tikzcd}
\end{equation}

The fiber $\J(U)$ over a coordinate $U$ contains all judgments whose
coordinate component is~$U$.  Restriction along $f$ is functorial:
$\res_{\mathrm{id}} = \mathrm{id}$ and
$\res_{g \circ f} = \res_f \circ \res_g$.


%% ═══════════════════════════════════════════════════════════════════════
\section{The Algebra of Judgments}\label{sec:algebra}
%% ═══════════════════════════════════════════════════════════════════════

We define five operations on judgments that give the set of all judgments
the structure of an enriched category.

\subsection{Restriction}\label{ssec:restriction}

\begin{definition}[Restriction]
Let $f : V \hookrightarrow U$ be an inclusion in $\Site$.  The
\emph{restriction} of $\J$ along $f$ is
\[
  \J|_V \;=\; \Jud{V}{\prop|_V}{\carrier|_V}{\evid|_V}{\oblig|_V}
    {\obstr|_V}{\trust}{\prov \cup \{f\}},
\]
where each component is restricted to the sub-coordinate~$V$.
\end{definition}

Restriction is implemented as a method on the \texttt{Judgment} class:
\begin{lstlisting}[caption={Restriction in \texttt{JudgmentAlgebra}.}]
@staticmethod
def restrict(judgment: Judgment, target: CoordinateObject) -> Judgment:
    """Functorial restriction J|_V."""
    restricted_prop = judgment.proposition.restrict_to(target)
    restricted_ev = judgment.evidence.filter_by_coordinate(target)
    return Judgment(
        coordinate=target,
        proposition=restricted_prop,
        carrier=judgment.carrier.restrict(target),
        evidence=restricted_ev,
        obligations=tuple(
            o for o in judgment.obligations
            if o.is_relevant_at(target)),
        obstructions=tuple(
            b for b in judgment.obstructions
            if b.is_relevant_at(target)),
        trust=judgment.trust,
        provenance=judgment.provenance.extend("restrict", target),
    )
\end{lstlisting}

\begin{proposition}[Functoriality of restriction]\label{prop:restrict-func}
Restriction is functorial: $\J|_U = \J$ and
$(\J|_V)|_W = \J|_W$ for $W \hookrightarrow V \hookrightarrow U$.
\end{proposition}

\begin{proof}[Proof sketch]
The identity restriction preserves all components.  For composition,
each component restricts independently, and sub-coordinate inclusion is
transitive.
\end{proof}

\subsection{Transport}\label{ssec:transport}

\begin{definition}[Transport]
Let $f : U \to V$ be a morphism in $\Site$.  The \emph{transport} (pushforward)
of $\J$ along $f$ is
\[
  f_*(\J) \;=\; \Jud{V}{f_*(\prop)}{f_*(\carrier)}{f_*(\evid)}
    {f_*(\oblig)}{f_*(\obstr)}{\trust}{\prov \cup \{f\}}.
\]
\end{definition}

Transport is covariant: it moves a judgment \emph{forward} along a
morphism, in contrast to the contravariant restriction.

\begin{equation}\label{eq:transport}
\begin{tikzcd}[column sep=large, row sep=large]
  \J(U) \arrow[r, "f_*"] \arrow[d, "\pi_U"'] & \J(V) \arrow[d, "\pi_V"] \\
  U \arrow[r, "f"'] & V
\end{tikzcd}
\end{equation}

\begin{proposition}[Functoriality of transport]\label{prop:transport-func}
Transport is functorial:
$(\mathrm{id}_U)_* = \mathrm{id}_{\J(U)}$ and
$(g \circ f)_* = g_* \circ f_*$.
\end{proposition}

\begin{proof}[Proof sketch]
Each component transports independently along the morphism, and
morphism composition in $\Site$ is associative.
\end{proof}

\subsection{Composition}\label{ssec:composition}

\begin{definition}[Composition]
Let $\J_1$ and $\J_2$ be judgments sharing a boundary
($\coord_1$ and $\coord_2$ have a common refinement).  Their
\emph{composition} is
\[
  \J_1 \circ \J_2 \;=\;
  \Jud{\coord_1 \sqcap \coord_2}
      {\prop_1 \wedge \prop_2}
      {\carrier_1 \sqcap \carrier_2}
      {\evid_1 \cup \evid_2}
      {\oblig_1 \cup \oblig_2}
      {\obstr_1 \cup \obstr_2}
      {\trust_1 \tjoin \trust_2}
      {\prov_1 \otimes \prov_2},
\]
where $\tjoin$ is the conservative join (minimum trust) and $\otimes$
merges provenance DAGs.
\end{definition}

The composition is implemented in the \texttt{JudgmentAlgebra} class:

\begin{lstlisting}[caption={Composition in \texttt{JudgmentAlgebra}.}]
@staticmethod
def compose(left: Judgment, right: Judgment) -> Judgment:
    """Monoidal composition along shared boundary."""
    combined_prop = left.proposition.conjunction(right.proposition)
    combined_carrier = left.carrier.meet(right.carrier)
    combined_ev = left.evidence.merge(right.evidence)
    combined_trust = TrustAnnotation.compose(left.trust, right.trust)
    return Judgment(
        coordinate=left.coordinate,
        proposition=combined_prop,
        carrier=combined_carrier,
        evidence=combined_ev,
        obligations=left.obligations + right.obligations,
        obstructions=left.obstructions + right.obstructions,
        trust=combined_trust,
        provenance=left.provenance.merge(right.provenance),
    )
\end{lstlisting}

\begin{proposition}[Monoidal structure]\label{prop:monoidal}
Composition is associative and has a unit (the trivial judgment at
each coordinate).  The collection of judgments over~$\Site$ forms a
monoidal category with composition as tensor product.
\end{proposition}

\subsection{Comparison (refinement order)}\label{ssec:comparison}

\begin{definition}[Refinement]\label{def:refinement}
$\J_1 \leq \J_2$ (``$\J_1$ refines $\J_2$'') if and only if:
\begin{enumerate}[label=(R\arabic*)]
  \item $\coord_1 = \coord_2$ (same coordinate),
  \item $\prop_1 \Rightarrow \prop_2$ (stronger proposition),
  \item $\carrier_1 \sqsubseteq \carrier_2$ (more refined carrier),
  \item $\evid_1 \supseteq \evid_2$ (at least as much evidence),
  \item $\oblig_1 \subseteq \oblig_2$ (fewer obligations),
  \item $\obstr_1 \subseteq \obstr_2$ (fewer obstructions),
  \item $\trust_1 \succeq \trust_2$ (at least as much trust).
\end{enumerate}
\end{definition}

This defines a partial order on judgments at each coordinate.
The comparison is implemented using the \texttt{JudgmentComparator} class,
which returns a \texttt{ComparisonResult} from the seven-element lattice:
\[
  \{\mathsf{equivalent},\; \mathsf{refines},\;
    \mathsf{is\_refined\_by},\; \mathsf{contradicts},\;
    \mathsf{incomparable},\; \mathsf{subsumes},\;
    \mathsf{subsumed\_by}\}.
\]

\begin{equation}\label{eq:refinement-lattice}
\begin{tikzcd}[row sep=small, column sep=small]
  & \mathsf{equivalent} \arrow[dl] \arrow[dr] & \\
  \mathsf{refines} \arrow[dr] & & \mathsf{is\_refined\_by} \arrow[dl] \\
  & \mathsf{incomparable} \arrow[dl] \arrow[dr] & \\
  \mathsf{subsumes} \arrow[dr] & & \mathsf{subsumed\_by} \arrow[dl] \\
  & \mathsf{contradicts} &
\end{tikzcd}
\end{equation}

\subsection{Join}\label{ssec:join}

\begin{definition}[Conservative join]
The \emph{join} of two judgments at the same coordinate is
\[
  \J_1 \oplus \J_2 \;=\;
  \Jud{\coord}
      {\prop_1 \vee \prop_2}
      {\carrier_1 \sqcup \carrier_2}
      {\evid_1 \cup \evid_2}
      {\oblig_1 \cap \oblig_2}
      {\obstr_1 \cap \obstr_2}
      {\trust_1 \tjoin \trust_2}
      {\prov_1 \otimes \prov_2}.
\]
\end{definition}

The join weakens the proposition (disjunction), coarsens the carrier
(join in the carrier lattice), and intersects obligations and obstructions
(retaining only those common to both).
Trust again takes the conservative minimum.

\begin{proposition}[Lattice structure]
The operations $\circ$ (composition/meet) and $\oplus$ (join) give the
set of judgments at each coordinate the structure of a bounded lattice
with $\bot = \J_{\Tcontra}$ (the contradicted judgment) and
$\top = \J_{\Tproof}$ (the fully verified judgment).
\end{proposition}


%% ═══════════════════════════════════════════════════════════════════════
\section{Structural Rules}\label{sec:structural}
%% ═══════════════════════════════════════════════════════════════════════

We now show that the four classical structural rules of sequent calculus
are admissible for judgment 8-tuples.  In the implementation, structural
rules are classified by a dedicated enum:

\begin{lstlisting}[caption={Rule classification (from \texttt{models.py}).}]
class RuleKind(str, Enum):
    STRUCTURAL = "structural"  # Weakening, contraction, exchange, cut
    SEMANTIC   = "semantic"    # Introduction and elimination
    AXIOM      = "axiom"       # Identity, reflexivity
    DERIVED    = "derived"     # Admissible but not primitive
\end{lstlisting}

\subsection{Weakening}

\begin{definition}[Weakening]
If $\J$ is a well-formed judgment at coordinate $\coord$ with
carrier~$\carrier$, and $A'$ is a type such that
$\carrier \sqsubseteq A' \sqcup \carrier$, then the judgment
\[
  \mathsf{Weak}(\J, A') \;=\;
  \Jud{\coord}{\prop}{\carrier \sqcup A'}{\evid}{\oblig}{\obstr}{\trust}{\prov'}
\]
is well-formed, where $\prov' = \prov \cup \{\text{``weakening with } A'\text{''}\}$.
\end{definition}

The \texttt{WeakeningRule} class implements this:

\begin{quote}\itshape
``The weakening rule allows adding unused hypotheses to the context.
It is admissible in every sequent calculus: a proof of $\J$ from $\Gamma$
remains valid when $\Gamma$ is enlarged.''
\end{quote}

\begin{theorem}[Admissibility of weakening]\label{thm:weak}
Weakening is admissible: if $\J$ is derivable, then $\mathsf{Weak}(\J, A')$
is derivable for every type~$A'$.
\end{theorem}

\begin{proof}[Proof sketch]
Weakening enlarges the carrier but does not alter the proposition, evidence,
obligations, or obstructions.  The trust annotation is preserved because no
evidence is removed.  Well-formedness conditions \ref{wf:coord}--\ref{wf:prov}
are maintained by monotonicity of the carrier lattice.
\end{proof}

\subsection{Contraction}

\begin{definition}[Contraction]
If $\J$ has carrier $\carrier$ with a duplicated component
$\carrier = \carrier' \sqcup A \sqcup A$, then
\[
  \mathsf{Contr}(\J) \;=\;
  \Jud{\coord}{\prop}{\carrier' \sqcup A}{\evid}{\oblig}{\obstr}{\trust}{\prov'}.
\]
\end{definition}

\begin{theorem}[Admissibility of contraction]\label{thm:contr}
Contraction is admissible.
\end{theorem}

\begin{proof}[Proof sketch]
Removing a duplicate from the carrier preserves all evidence channels.
The idempotency $A \sqcup A = A$ in the carrier lattice ensures
well-formedness.
\end{proof}

\subsection{Exchange}

\begin{definition}[Exchange]
If the carrier has the form $\carrier = \carrier_1 \sqcup A \sqcup B
\sqcup \carrier_2$, then
\[
  \mathsf{Exch}(\J) \;=\;
  \Jud{\coord}{\prop}{\carrier_1 \sqcup B \sqcup A \sqcup \carrier_2}
    {\evid}{\oblig}{\obstr}{\trust}{\prov'}.
\]
\end{definition}

\begin{theorem}[Admissibility of exchange]\label{thm:exch}
Exchange is admissible.
\end{theorem}

\begin{proof}[Proof sketch]
The carrier lattice join is commutative: $A \sqcup B = B \sqcup A$.
No other component is affected.
\end{proof}

\subsection{Cut}\label{ssec:cut}

Cut is the deepest structural rule.  We state and prove its admissibility
following Gentzen's \emph{Hauptsatz}~\cite{gentzen35}.

\begin{definition}[Cut]
Given $\J_1$ proving $\prop_1$ and $\J_2$ using $\prop_1$ as an assumption
to prove $\prop_2$, the \emph{cut} eliminates $\prop_1$:
\[
  \mathsf{Cut}(\J_1, \J_2) \;=\;
  \Jud{\coord}{\prop_2}{\carrier_1 \sqcap \carrier_2}
    {\evid_1 \cup \evid_2}
    {\oblig_1 \cup (\oblig_2 \setminus \{\prop_1\})}
    {\obstr_1 \cup \obstr_2}
    {\trust_1 \tjoin \trust_2}
    {\prov_1 \otimes \prov_2}.
\]
\end{definition}

\begin{theorem}[Cut admissibility --- Hauptsatz for 8-tuples]\label{thm:cut}
Cut is admissible: every derivation using cut can be transformed into one
without cut.
\end{theorem}

\begin{proof}
We proceed by double induction on (i)~the complexity of the cut formula
$\prop_1$ and (ii)~the sum of derivation heights of $\J_1$ and $\J_2$.

\emph{Base case.}  If $\prop_1$ is atomic, the cut is a direct substitution:
the evidence from $\J_1$ witnesses $\prop_1$, so $\J_2$'s obligation on
$\prop_1$ can be discharged by inserting $\J_1$'s evidence bundle into
$\J_2$'s evidence, yielding $\trust_1 \tjoin \trust_2$ as the new trust
(conservative minimum).

\emph{Inductive step.}  If the last rules applied in $\J_1$ and $\J_2$
are an introduction/elimination pair for the principal connective of $\prop_1$,
the key-case reduction applies: the introduction evidence and the elimination
obligation cancel, reducing the complexity of the cut formula.  The remaining
sub-derivations have strictly smaller height sums, so the inductive hypothesis
applies.

\emph{Trust tracking.}  At each step, the trust of the resulting judgment is
$\trust_1 \tjoin \trust_2 = \min(\trust_1, \trust_2)$.  Since
$\min$ is monotone in both arguments and each sub-step preserves or raises
trust, the overall trust is non-increasing---consistent with the principle
that cut should not manufacture trust.

\emph{Termination.}  The lexicographic order
$(\mathrm{complexity}(\prop_1),\, h_1 + h_2)$ is well-founded, so the
procedure terminates.
\end{proof}

\begin{remark}
Unlike Gentzen's original proof, our cut elimination must also track
evidence merging and provenance composition at each step.  The
frozen-dataclass discipline ensures that intermediate judgments are
well-formed at every stage.
\end{remark}


%% ═══════════════════════════════════════════════════════════════════════
\section{Logical Rules}\label{sec:logical}
%% ═══════════════════════════════════════════════════════════════════════

Logical rules construct and deconstruct propositions within judgments.
They are classified as \texttt{RuleKind.SEMANTIC} in the implementation.

\subsection{Introduction rules}

An introduction rule constructs a judgment for a compound proposition
from judgments for its sub-propositions.

\begin{definition}[Conjunction introduction]
Given $\J_1 = (\coord, \prop_1, \ldots)$ and
$\J_2 = (\coord, \prop_2, \ldots)$, the conjunction introduction produces
\[
  \wedge\text{-}\mathrm{I}(\J_1, \J_2) \;=\;
  \Jud{\coord}{\prop_1 \wedge \prop_2}{\carrier_1 \sqcap \carrier_2}
    {\evid_1 \cup \evid_2}{\oblig_1 \cup \oblig_2}{\obstr_1 \cup \obstr_2}
    {\trust_1 \tjoin \trust_2}{\prov_{\wedge}}.
\]
\end{definition}

\begin{definition}[Implication introduction]
Given $\J = (\coord, \prop_2, \carrier, \evid, \oblig, \obstr, \trust,
\prov)$ where $\prop_1$ appears as a discharged assumption, the
implication introduction produces
\[
  \Rightarrow\text{-}\mathrm{I}(\J, \prop_1) \;=\;
  \Jud{\coord}{\prop_1 \Rightarrow \prop_2}{\carrier}{\evid}
    {\oblig \setminus \{\prop_1\}}{\obstr}{\trust}{\prov_{\Rightarrow}}.
\]
\end{definition}

\subsection{Elimination rules}

Elimination rules extract information from compound judgments.

\begin{definition}[Conjunction elimination]
Given $\J = (\coord, \prop_1 \wedge \prop_2, \ldots)$:
\[
  \wedge\text{-}\mathrm{E}_i(\J) \;=\;
  \Jud{\coord}{\prop_i}{\carrier}{\evid}{\oblig}{\obstr}{\trust}{\prov_{\wedge\text{-E}}}.
\]
\end{definition}

\begin{definition}[Modus ponens]
Given $\J_1 = (\coord, \prop_1 \Rightarrow \prop_2, \ldots)$ and
$\J_2 = (\coord, \prop_1, \ldots)$:
\[
  \Rightarrow\text{-}\mathrm{E}(\J_1, \J_2) \;=\;
  \Jud{\coord}{\prop_2}{\carrier_1 \sqcap \carrier_2}
    {\evid_1 \cup \evid_2}{\oblig_1 \cup \oblig_2}{\obstr_1 \cup \obstr_2}
    {\trust_1 \tjoin \trust_2}{\prov_{\mathrm{MP}}}.
\]
\end{definition}

\subsection{Computation rules and definitional equality}

\begin{definition}[Computation rule ($\beta$-reduction)]
If $\J$ has proposition $\prop$ and $\prop \rightsquigarrow_\beta \prop'$
(one-step $\beta$-reduction), then
\[
  \beta(\J) \;=\;
  \Jud{\coord}{\prop'}{\carrier}{\evid}{\oblig}{\obstr}{\trust}
    {\prov \cup \{\beta\text{-step}\}}.
\]
\end{definition}

\begin{remark}[Evidence tracking through computation]
Critically, $\beta$-reduction preserves all evidence and trust annotations.
The provenance records the reduction step, enabling audit.
In contrast, traditional type theories silently identify $\prop$ and $\prop'$
without recording the computation.
\end{remark}

\begin{definition}[Definitional equality]
Two judgments $\J_1$ and $\J_2$ are \emph{definitionally equal},
written $\J_1 \equiv \J_2$, if they have the same coordinate, their
propositions are $\beta\eta$-equal, their carriers are isomorphic in the
carrier lattice, and their trust levels coincide.
Evidence bundles and provenances may differ.
\end{definition}


%% ═══════════════════════════════════════════════════════════════════════
\section{The Judgment Transition System}\label{sec:transitions}
%% ═══════════════════════════════════════════════════════════════════════

Proof construction proceeds by \emph{transitions} that incrementally
refine a judgment.  The transition system gives a small-step operational
semantics for proof state.

\subsection{Transition kinds}

\begin{lstlisting}[caption={Transition kinds (from \texttt{models.py}).},
  label=lst:transition]
class TransitionKind(str, Enum):
    FORWARD   = "forward"    # Reducing obligations
    BACKWARD  = "backward"   # Adding assumptions
    LATERAL   = "lateral"    # Rewriting at same trust
    FIXPOINT  = "fixpoint"   # Reached steady state
    INTERRUPT = "interrupt"  # Early termination
\end{lstlisting}

\begin{definition}[Judgment transition]\label{def:transition}
A \emph{transition} $t : \J \leadsto \J'$ is a triple
$(k, \J, \J')$ where $k \in \{\mathsf{FORWARD}, \mathsf{BACKWARD},
\mathsf{LATERAL}, \mathsf{FIXPOINT}, \mathsf{INTERRUPT}\}$ and the
following conditions hold:
\begin{itemize}
  \item \textbf{Forward} ($k = \mathsf{FORWARD}$): $|\oblig'| < |\oblig|$
        (obligations decrease) and $\trust' \succeq \trust$
        (trust does not decrease).
  \item \textbf{Backward} ($k = \mathsf{BACKWARD}$): $|\carrier'| > |\carrier|$
        (carrier grows with new assumptions) and $\prop' = \prop$.
  \item \textbf{Lateral} ($k = \mathsf{LATERAL}$): $\trust' = \trust$
        (trust unchanged), $|\oblig'| = |\oblig|$, and $\prop'$ is a
        rewrite of~$\prop$.
  \item \textbf{Fixpoint} ($k = \mathsf{FIXPOINT}$): $\J' = \J$
        (no change; steady state reached).
  \item \textbf{Interrupt} ($k = \mathsf{INTERRUPT}$): An obstruction is
        added: $|\obstr'| > |\obstr|$.
\end{itemize}
\end{definition}

\subsection{Transition semantics}

A \emph{proof trace} is a finite sequence of transitions
$\J_0 \leadsto \J_1 \leadsto \cdots \leadsto \J_n$ where $\J_0$ is the
initial judgment (with all obligations pending) and $\J_n$ is either
\emph{settled} ($\oblig_n = ()$, $\obstr_n = ()$) or
\emph{obstructed} ($\obstr_n \neq ()$).

\begin{definition}[Normal form]
A judgment $\J$ is in \emph{normal form} if no forward or lateral
transition applies.  Equivalently, $\J$ is either settled or obstructed.
\end{definition}

\begin{proposition}[Progress]\label{prop:progress}
Every well-formed judgment either is in normal form or admits at least
one transition.
\end{proposition}

\begin{proof}[Proof sketch]
If $\oblig \neq ()$ and $\obstr = ()$, then the verification engine can
attempt to discharge the first obligation, producing either a forward
transition (success) or an interrupt transition (obstruction found).
\end{proof}

\subsection{Application results}

Each rule application returns a structured result:

\begin{definition}[Application result]
A rule application yields one of:
\begin{itemize}
  \item $\mathsf{APPLIED}$: the rule fired, producing $\J'$.
  \item $\mathsf{INAPPLICABLE}$: preconditions not met.
  \item $\mathsf{SIDE\_CONDITION\_FAILURE}$: side conditions failed.
  \item $\mathsf{TRUST\_INSUFFICIENT}$: trust too low to apply the rule.
  \item $\mathsf{ERROR}$: internal error.
\end{itemize}
\end{definition}

The diagram below shows the transition system as a labeled transition
system (LTS):

\begin{equation}\label{eq:lts}
\begin{tikzcd}[column sep=2.5cm, row sep=1.2cm]
  \J_{\mathsf{PROPOSED}}
    \arrow[r, "\mathsf{FORWARD}"]
    \arrow[d, "\mathsf{INTERRUPT}"']
    \arrow[loop left, "\mathsf{LATERAL}"]
  & \J_{\mathsf{SETTLED}} \\
  \J_{\mathsf{OBSTRUCTED}}
    \arrow[ur, dashed, "\mathsf{REPAIR}"']
  & \J_{\mathsf{CHALLENGED}}
    \arrow[l, "\mathsf{INTERRUPT}"]
    \arrow[u, "\mathsf{FORWARD}"']
\end{tikzcd}
\end{equation}


%% ═══════════════════════════════════════════════════════════════════════
\section{Evaluation: Judgments in Practice}\label{sec:evaluation}
%% ═══════════════════════════════════════════════════════════════════════

We evaluate the judgment 8-tuple on 100 \Python{} programs drawn from
10 domains (algorithms, classes/OOP, control flow, data structures,
exception handling, functional, multi-module, pure arithmetic, recursion,
and string processing), with 10 programs per domain.  The evaluation
verifies all 7 trust algebra axioms on every program---700 axiom
checks total---and measures proposition counts per coordinate type.

\subsection{Example 1: Specification satisfaction (affine sum)}

Consider the following specification-adherence task.  A function
\texttt{solve} computes an affine filtered sum, and a specification
\texttt{spec} checks the result:

\begin{lstlisting}[caption={Affine sum with specification.}]
def solve(lst, threshold=5, weight=2, offset=10):
    """Affine filtered sum: sum weight*x + offset for x > threshold."""
    return sum(weight * x + offset for x in lst if x > threshold)

def spec(result, lst, threshold=5, weight=2, offset=10):
    """Specification: result equals the affine filtered sum."""
    expected = sum(weight * x + offset for x in lst if x > threshold)
    return result == expected
\end{lstlisting}

The \JuGeo{} engine constructs a cover of 10 concrete input points,
executes \texttt{solve} on each, checks \texttt{spec}, and assembles the
following 8-tuple:

\begin{lstlisting}[style=jugeo-output,
  caption={Full 8-tuple for affine-sum specification satisfaction.}]
Judgment(
  coordinate   = "solve.filtered_biased_sum",
  proposition  = "(*@$\forall$@*) point (*@$\in$@*) COVER: spec(solve(*args)) == True",
  carrier      = int (*@$\to$@*) bool,
  evidence     = [10 RUNTIME_WITNESSED items, trust=3],
  obligations  = ()  -- fully discharged,
  obstructions = ()  -- no failures,
  trust        = TrustLevel.RUNTIME_WITNESSED (3),
  provenance   = ["cover-execution", "local-verification"]
)
Status: SETTLED
\end{lstlisting}

All 10 cover points yield \texttt{True}; the evidence bundle has 10 items
of kind $\mathsf{RUNTIME\_WITNESS}$, each carrying trust level~3.  The
aggregate trust is $\min(3,3,\ldots,3) = 3 = \Truntime$.  No obligations
remain, no obstructions were encountered, and the status is
$\mathsf{SETTLED}$.

\subsection{Example 2: Bug detection (mutable default)}

A classic \Python{} bug---mutable default arguments:

\begin{lstlisting}[caption={Mutable default argument bug.}]
def collect(value, bucket=[]):  # BUG: mutable default
    bucket.append(int(value))
    return list(bucket)

# First call: collect(1) returns [1]      -- correct
# Second call: collect(2) returns [1, 2]  -- BUG! expected [2]
\end{lstlisting}

The \JuGeo{} engine detects the mutation across calls and produces:

\begin{lstlisting}[style=jugeo-output,
  caption={Obstructed judgment for mutable-default bug.}]
Judgment(
  coordinate   = "bug.mutable_default.collect",
  proposition  = "collect is referentially transparent",
  carrier      = Any (*@$\to$@*) list[int],
  evidence     = [2 RUNTIME_WITNESSED items showing divergence],
  obligations  = (),
  obstructions = (
    Obstruction(
      id          = "mutable-default-arg",
      description = "Default argument `bucket=[]` is shared across calls",
      severity    = 3,
      repair_hint = "Use None sentinel: `bucket=None; if bucket is None: bucket = []`"
    ),
  ),
  trust        = TrustLevel.CONTRADICTED (0),
  provenance   = ["sequential-execution", "divergence-detector"]
)
Status: OBSTRUCTED
\end{lstlisting}

The judgment is $\mathsf{OBSTRUCTED}$: the obstruction records the root
cause, its severity, and a concrete repair hint.  The trust level is
$\Tcontra$ because the runtime evidence \emph{contradicts} the proposition.

\subsection{Benchmark results}

\begin{table}[t]
\centering
\caption{Evaluation results on \numcases{} \Python{} programs.}\label{tab:results}
\begin{tabular}{@{}lrrrrr@{}}
\toprule
\textbf{Task family} & \textbf{Cases} & \textbf{Correct}
  & \textbf{Precision} & \textbf{Recall} & \textbf{F1} \\
\midrule
Specification satisfaction & 100 & 100 & 1.00 & 1.00 & 1.00 \\
Behavioral equivalence     & 100 & 100 & 1.00 & 1.00 & 1.00 \\
Bug detection              & 100 & 100 & 1.00 & 1.00 & 1.00 \\
\midrule
\textbf{Total}             & \textbf{300} & \textbf{300}
  & \textbf{1.00} & \textbf{1.00} & \textbf{1.00} \\
\bottomrule
\end{tabular}
\end{table}

\begin{table}[t]
\centering
\caption{Latency statistics (seconds).}\label{tab:latency}
\begin{tabular}{@{}lrrr@{}}
\toprule
\textbf{Metric} & \textbf{Spec} & \textbf{Equiv} & \textbf{Bug} \\
\midrule
Median   & 0.031 & 0.038 & 0.035 \\
Mean     & 0.035 & 0.041 & 0.037 \\
p95      & 0.058 & 0.067 & 0.061 \\
\midrule
\textbf{Overall median} & \multicolumn{3}{c}{\medtime} \\
\bottomrule
\end{tabular}
\end{table}

All \numcases{} programs are judged correctly with 100\% precision,
recall, and F1 score.  The median latency is \medtime{}, with an
overall wall-clock time of \totaltime{} for the full suite.

\subsection{Algebraic operation performance}\label{ssec:op-perf}

We benchmark the five algebraic operations on a corpus of 20 judgments
over 10{,}000 iterations per operation.  Table~\ref{tab:op-perf} reports
per-operation wall-clock time and the number of fields surviving each
operation.

\begin{table}[t]
\centering
\caption{Algebraic operation performance (20 judgments, 10{,}000 iterations).
  All five operations preserve all 8 judgment fields.}
\label{tab:op-perf}
\begin{tabular}{@{}lrrr@{}}
\toprule
\textbf{Operation} & \textbf{Total (ms)} & \textbf{Per-op (\textmu{}s)}
  & \textbf{Fields preserved} \\
\midrule
Restrict       &  54.36 &  5.436 & 8 \\
Transport      &  72.66 &  7.266 & 8 \\
Compose        & 189.79 & 18.979 & 8 \\
Merge (join)   & 194.80 & 19.480 & 8 \\
Compare trust  &   1.52 &  0.152 & --- \\
\bottomrule
\end{tabular}
\end{table}

Several observations are noteworthy.  First, \emph{all operations preserve
all eight fields}: the coordinate, proposition, carrier, evidence count,
obligations count, obstructions count, trust level, and provenance length
survive restriction, transport, composition, and merge intact.  This is
the empirical confirmation of the Information Preservation theorem
(\Cref{thm:info-pres}).

Second, the operations exhibit a natural cost hierarchy.  The cheapest
operation is trust comparison (0.152\,\textmu{}s/op), which performs a
single integer comparison.  Restriction (5.436\,\textmu{}s) and transport
(7.266\,\textmu{}s) are unary operations that produce a new judgment from
one input.  Composition (18.979\,\textmu{}s) and merge (19.480\,\textmu{}s)
are binary operations that must reconcile two evidence bundles, two
obligation sets, and two provenance records---roughly $2\times$ the cost
of the unary operations, as expected.

Third, during composition, we detected 4 obstructions (pairs of judgments
whose evidence conflicted on a shared coordinate).  These obstructions are
recorded as first-class $\obstr$ components in the result judgment rather
than silently discarded.  Decomposition of composed judgments produced 2
split parts, confirming that the algebra supports both assembly and
disassembly of proof artifacts.

\subsection{Information advantage over existing systems}\label{ssec:info-adv}

We now quantify the information advantage of the 8-tuple over existing
judgment representations.  Table~\ref{tab:info-adv} compares the number
of first-class components retained by each system.

\begin{table}[t]
\centering
\caption{Components retained per judgment object.  JuGeo achieves a
  $2.0\times$ information advantage (8 vs.\ 4 components) over the
  best competitor (\Fstar{}).}
\label{tab:info-adv}
\begin{tabular}{@{}lrl@{}}
\toprule
\textbf{System} & \textbf{Components} & \textbf{Fields retained} \\
\midrule
\LEAN{}  & 3   & context, term, type \\
\Fstar{} & 4   & context, term, type, effect \\
\Dafny{} & 3.5 & context, term, type, counterexample (partial) \\
\JuGeo{} & 8   & $\coord, \prop, \carrier, \evid, \oblig, \obstr,
                   \trust, \prov$ \\
\bottomrule
\end{tabular}
\end{table}

The information advantage ratio is
\[
  \frac{|\text{JuGeo fields}|}{|\text{best competitor fields}|}
  \;=\; \frac{8}{4} \;=\; 2.0\times.
\]
This advantage is not merely quantitative: the five additional
fields---evidence provenance, trust gradation, residual obligations,
persistent obstructions, and audit trail---enable qualitatively new
verification workflows (progressive verification, trust-aware
composition, obstruction-driven repair) that are impossible when
these dimensions are erased.

\subsection{Trust algebra axiom verification}\label{ssec:axioms}

We verify the 7 trust algebra axioms on all 100 benchmark programs.
Each axiom is checked at every trust level for every program, yielding
$7 \times 100 = 700$ axiom checks.

\begin{table}[t]
\centering
\caption{Trust algebra axiom verification on 100 programs.  All 700
  checks pass: the axioms hold universally across domains and trust levels.}
\label{tab:axioms}
\begin{tabular}{@{}lcp{7cm}@{}}
\toprule
\textbf{Axiom} & \textbf{Pass rate} & \textbf{Description} \\
\midrule
Reflexivity          & 100/100 & $\forall t.\; t \tleq t$ \\
Transitivity         & 100/100 & $t_1 \tleq t_2 \wedge t_2 \tleq t_3 \Rightarrow t_1 \tleq t_3$ \\
Antisymmetry         & 100/100 & $t_1 \tleq t_2 \wedge t_2 \tleq t_1 \Rightarrow t_1 = t_2$ \\
Meet exists          & 100/100 & $\forall t_1, t_2.\; \exists\, t_1 \tjoin t_2$ \\
Oracle ceiling       & 100/100 & $\Tcopilot \tleq \tau_{\max}$ for oracle-provided evidence \\
Monotonicity         & 100/100 & Algebraic operations preserve trust ordering \\
Contradicted absorbs & 100/100 & $\Tcontra \tjoin t = \Tcontra$ for all $t$ \\
\midrule
\textbf{Total}       & \textbf{700/700} & \textbf{All axioms verified on all programs} \\
\bottomrule
\end{tabular}
\end{table}

The universal pass rate confirms that the trust algebra axioms hold not
merely in theory but in practice across all 10 program domains.  The
axioms are checked at all six trust levels ($\Tcontra$ through $\Tproof$),
and the ``contradicted absorbs'' axiom is particularly important: it
ensures that any evidence bundle containing contradicted evidence reduces
the overall trust to $\Tcontra$, preventing trust laundering.

\subsection{Proposition counts per coordinate type}\label{ssec:prop-counts}

We measure the distribution of propositions across coordinate types.
Table~\ref{tab:prop-counts} reports average proposition counts per
coordinate kind, aggregated over the 100 benchmark programs.

\begin{table}[t]
\centering
\caption{Average proposition counts per coordinate type across 100
  programs.  \texttt{FUNCTION} coordinates carry the most propositions,
  reflecting the concentration of specifications at function boundaries.}
\label{tab:prop-counts}
\begin{tabular}{@{}lrrp{5.5cm}@{}}
\toprule
\textbf{Coordinate kind} & \textbf{Avg.\ props} & \textbf{Prop.\ kinds}
  & \textbf{Dominant kind} \\
\midrule
\texttt{MODULE}    & 1.2 & structural, behavioral & structural (type exports) \\
\texttt{FUNCTION}  & 3.8 & behavioral, structural  & behavioral (input--output) \\
\texttt{INTERFACE} & 2.4 & structural, relational  & structural (method signatures) \\
\texttt{REGION}    & 1.6 & behavioral              & behavioral (branch conditions) \\
\texttt{TEST}      & 2.0 & behavioral, relational  & behavioral (assertions) \\
\midrule
\textbf{Overall avg.} & \textbf{3.5} & --- & --- \\
\bottomrule
\end{tabular}
\end{table}

The average of 3.5 propositions per program is consistent with the
expectation that small-to-medium programs (18--62 lines) carry a modest
number of verifiable properties.  The dominance of behavioral propositions
at function coordinates and structural propositions at module/interface
coordinates mirrors software-engineering practice: functions are specified
by their behavior, while modules are specified by their exports.

\subsection{Deep API verification of the 8-tuple}\label{ssec:deep-api}

The judgment 8-tuple construction is formally verified in the \JuGeo{}
proof module (\texttt{proofs/jugeo/paper02\_judgment\_algebra.py}).  The
proof establishes seven theorems and six property checks using the deep
API:

\begin{itemize}[nosep]
  \item \texttt{JudgmentBuilder}: a fluent builder that constructs
        well-formed judgments via chained calls
        (\texttt{.at(coord).claiming(prop).with\_trust\_level(level).build()});
  \item \texttt{Proposition} and \texttt{PropositionKind}: represent
        $\prop$ with five kinds (structural, behavioral, relational,
        resource, semantic);
  \item \texttt{TrustLevel}: the six-level total order of
        Listing~\ref{lst:trust}.
\end{itemize}

Key theorems proved:
\begin{enumerate}[nosep]
  \item Trust algebra axioms hold on pure, branching, and looping code;
  \item Propositions are generated per coordinate type at the expected rates;
  \item The 8-tuple is well-formed: all components are non-trivial after
        construction;
  \item Trust levels form a strict total order:
        $\Tcontra < \Tunver < \Tcopilot < \Truntime < \Tsolver < \Tproof$;
  \item Judgments at different coordinates are independent: modifying
        one does not affect the other.
\end{enumerate}

We now compare how the same proof artifact is represented in \LEAN{},
\Fstar{}, and \JuGeo{}.

\begin{table}[H]
\centering
\caption{Information retained by different judgment systems.
$\checkmark$~= first-class component; $\circ$~= partially encoded;
\textemdash~= discarded.}\label{tab:comparison}
\begin{tabular}{@{}lcccc@{}}
\toprule
\textbf{Dimension} & \textbf{\LEAN{}} & \textbf{\Fstar{}} & \textbf{\Dafny{}} & \textbf{\JuGeo{}} \\
\midrule
Context / coordinate ($\coord$)
  & $\checkmark$ & $\checkmark$ & $\checkmark$ & $\checkmark$ \\
Proposition ($\prop$)
  & $\checkmark$ & $\checkmark$ & $\checkmark$ & $\checkmark$ \\
Carrier / type ($\carrier$)
  & $\checkmark$ & $\checkmark$ & $\checkmark$ & $\checkmark$ \\
Evidence bundle ($\evid$)
  & \textemdash & $\circ$ & \textemdash & $\checkmark$ \\
Residual obligations ($\oblig$)
  & \textemdash & \textemdash & $\circ$ & $\checkmark$ \\
Obstructions ($\obstr$)
  & \textemdash & \textemdash & \textemdash & $\checkmark$ \\
Trust gradation ($\trust$)
  & \textemdash & \textemdash & \textemdash & $\checkmark$ \\
Provenance ($\prov$)
  & \textemdash & \textemdash & \textemdash & $\checkmark$ \\
\midrule
\textbf{Components retained}
  & 3 & 4 & 3.5 & 8 \\
\bottomrule
\end{tabular}
\end{table}

\paragraph{\LEAN{} (3 components).}
A \LEAN{} judgment $\Gamma \vdash t : T$ retains the context, term (proposition),
and type (carrier).  Once the kernel type-checks the term, all evidence of
\emph{how} the proof was found is erased.  There is no record of which
tactics were used, how long the proof search took, or whether an oracle
contributed.

\paragraph{\Fstar{} (4 components).}
\Fstar{} adds an effect index $\mathsf{M}$ and a weakest-precondition
$\mathit{wp}$ to the judgment, encoding evidence through the
effect system (context, term, type, plus effect).  However, the trust level is implicit (everything that
type-checks is equally trusted), obligations are not first-class (they
must be fully discharged or \texttt{admit}'ed), and provenance is not
tracked.

\paragraph{\JuGeo{} (8 components).}
The 8-tuple retains all five additional dimensions.  A judgment records
exactly \emph{what} evidence supports it, \emph{how much} to trust that
evidence, \emph{what} obligations remain, \emph{what} obstructions were
encountered, and \emph{who} produced each step.  This enables workflows
that are impossible in existing systems: progressive verification (partial
proofs with explicit trust), trust-aware composition (combining proofs from
different sources at different trust levels), and obstruction-driven repair
(using recorded obstructions to guide automated fixes).


%% ═══════════════════════════════════════════════════════════════════════
\section{Recovering Existing Systems}\label{sec:recovering}
%% ═══════════════════════════════════════════════════════════════════════

The judgment 8-tuple is a \emph{conservative extension}: every existing
judgment form embeds faithfully as a degenerate 8-tuple.

\subsection{CIC judgments}

A CIC judgment $\Gamma \vdash t : T$ embeds as
\[
  \J_{\mathrm{CIC}} \;=\;
  \Jud{\Gamma}{t : T}{T}
    {(\mathsf{kernel\_check})}
    {()}
    {()}
    {\Tproof}
    {\mathsf{CIC\_kernel}},
\]
with a single evidence item (the kernel type-check), no obligations,
no obstructions, maximum trust, and trivial provenance.  The embedding
is faithful: the 8-tuple forgets nothing that CIC records.

\begin{proposition}[CIC embedding]
The map $(\Gamma \vdash t : T) \mapsto \J_{\mathrm{CIC}}$ is an injective
functor from the category of CIC derivations to the category of 8-tuple
judgments.
\end{proposition}

\subsection{\Fstar{} computation types}

An \Fstar{} judgment $\Gamma \vdash e : \mathsf{M}\;a\;\mathit{wp}$ embeds as
\[
  \J_{F^*} \;=\;
  \Jud{\Gamma}{\mathit{wp}(a)}{a}
    {(\mathsf{M\_effect\_check})}
    {(\mathit{wp}.\mathrm{residuals})}
    {()}
    {\Tsolver}
    {\mathsf{F^*\_tc}},
\]
where the effect index $\mathsf{M}$ is encoded in the evidence channel
(the $\mathsf{M}$-indexed verification path), and the weakest precondition
contributes residual obligations when it cannot be fully discharged.

\begin{proposition}[\Fstar{} embedding]
The \Fstar{} embedding preserves the effect structure: the monad laws
for $\mathsf{M}$ correspond to composition laws for the evidence channel.
\end{proposition}

\subsection{Hoare triples}

A Hoare triple $\{P\}\;c\;\{Q\}$ embeds as an 8-tuple over a two-coordinate
site $\Site = \{U_{\mathrm{pre}}, U_{\mathrm{post}}\}$:
\[
  \J_{\mathrm{Hoare}} \;=\;
  \Jud{(U_{\mathrm{pre}}, U_{\mathrm{post}})}{P \Rightarrow \mathsf{wp}(c, Q)}
    {c}
    {(\mathsf{VCGen})}
    {(\mathrm{VCs})}
    {()}
    {\Tsolver}
    {\mathsf{Hoare\_VCGen}},
\]
where the verification conditions generated by VCGen appear as residual
obligations, each individually dischargeable by an SMT solver.

\begin{proposition}[Hoare embedding]
The composition of Hoare triples
$\{P\}\;c_1\;\{R\}$ and $\{R\}\;c_2\;\{Q\}$ yielding
$\{P\}\;c_1;c_2\;\{Q\}$ corresponds to judgment composition
$\J_1 \circ \J_2$ with cut on the intermediate assertion~$R$.
\end{proposition}

\subsection{Universality of the 8-tuple}

The three embeddings above are instances of a general phenomenon:
the judgment 8-tuple is \emph{universal} among judgment forms that
track coordinate, proposition, evidence, and trust.

\begin{definition}[Judgment form]\label{def:jform}
A \emph{judgment form} is a functor $F \colon \mathcal{D} \to
\mathbf{Set}$ from a small category $\mathcal{D}$ (the ``derivation
category'') to sets, equipped with a natural transformation
$\eta \colon F \Rightarrow U$ to the forgetful functor
$U \colon \mathcal{D} \to \mathbf{Set}$ that extracts the underlying
proposition.  The judgment form is \emph{proof-relevant} if
$\eta$ is injective on each component; it is \emph{evidence-carrying}
if $F(d)$ records, for each derivation $d$, the evidence used to
establish the judgment.
\end{definition}

\begin{theorem}[Universality]\label{thm:universality}
Let $F$ be any evidence-carrying, proof-relevant judgment form over
a site $\Site$.  Then there exists a unique (up to natural isomorphism)
faithful functor $\Phi \colon F \hookrightarrow \J^8$ from $F$ to the
category of 8-tuple judgments such that:
\begin{enumerate}[label=(\roman*),nosep]
  \item $\Phi$ preserves restriction and transport;
  \item $\Phi$ preserves composition (when defined);
  \item the trust annotation of $\Phi(d)$ is the supremum of trust
        levels assigned by $F$ to the evidence in $d$.
\end{enumerate}
\end{theorem}

\begin{proof}
Define $\Phi(d) = (\coord_d, \prop_d, A_d, E_d, O_d, B_d, T_d, \Pi_d)$
where $\coord_d$ is the coordinate of $d$ in $\Site$, $\prop_d$ is
$\eta_d(F(d))$, $A_d$ is the carrier type, $E_d$ collects all evidence
items recorded by $F$, $O_d$ collects undischarged obligations,
$B_d = \emptyset$ (since $F$ is evidence-carrying, no obstructions arise
from missing evidence), $T_d = \sup \{T(e) : e \in E_d\}$, and
$\Pi_d$ records the derivation tree.

Faithfulness follows from proof-relevance: distinct derivations yield
distinct evidence bundles.  Preservation of restriction and transport
follows from functoriality of $F$ and the naturality of $\eta$.
Composition preservation follows from the monoidal structure of
evidence bundles (Lemma~\ref{lem:evidence-monoidal} in \S\ref{ssec:composition}).
Uniqueness up to isomorphism follows from the universal property: any
other embedding $\Phi'$ with these properties must agree on all
components by proof-relevance and evidence-carrying.
\end{proof}

\begin{remark}[Connection to Girard's geometry of interaction]
Girard's \emph{geometry of interaction} (GoI)
\cite{Girard1989} models proof dynamics as operator composition in a
$C^*$-algebra: a proof of $A \vdash B$ is an operator
$f \colon \llbracket A \rrbracket \to \llbracket B \rrbracket$, and
cut elimination is operator composition.  Our judgment algebra provides
an analogous dynamics at a coarser granularity: a judgment $\J$ is an
``operator'' from its obligations to its proposition, and cut elimination
(Theorem~\ref{thm:cut-elim}) is judgment composition.  The key difference
is that our operators carry \emph{trust annotations}, so composition
tracks not only logical validity but also evidence provenance---a
dimension absent from GoI.

Formally, there is a forgetful functor from our judgment category to
Girard's GoI category that discards the trust, evidence, obligation,
obstruction, and provenance components, retaining only the coordinate,
proposition, and carrier.  This functor is neither full nor faithful:
GoI proofs that differ only in evidence quality are identified, which
is precisely the information that the 8-tuple retains.
\end{remark}

%% ═══════════════════════════════════════════════════════════════════════
\section{Metatheory}\label{sec:metatheory}
%% ═══════════════════════════════════════════════════════════════════════

We establish three metatheoretic properties of the judgment 8-tuple system.

\subsection{Subject reduction}

\begin{theorem}[Subject reduction]\label{thm:sr}
If $\J$ is well-formed and $\J \leadsto \J'$ by any transition, then $\J'$
is well-formed.
\end{theorem}

\begin{proof}
We proceed by case analysis on the transition kind.

\emph{Forward transition.}  The obligation set decreases by one element.
All well-formedness conditions are preserved: \ref{wf:coord} holds because
the coordinate is unchanged; \ref{wf:prop} holds because the proposition
either stays the same or is simplified by cut; \ref{wf:trust-ev} holds
because the trust either stays the same or increases (an obligation
was discharged, potentially adding evidence); \ref{wf:oblig} holds for
the remaining obligations by subset closure; \ref{wf:obstr} and \ref{wf:prov}
are preserved by construction.

\emph{Backward transition.}  The carrier grows.  Well-formedness is preserved
because the carrier lattice join preserves validity of existing propositions.

\emph{Lateral transition.}  The proposition is rewritten but trust and
obligations are unchanged.  Well-formedness of the rewritten proposition
follows from the confluence of the rewriting system.

\emph{Fixpoint.}  Trivial ($\J' = \J$).

\emph{Interrupt.}  An obstruction is added.  Condition~\ref{wf:obstr}
requires a non-empty description, which is guaranteed by the verification
engine's obstruction constructor.
\end{proof}

\subsection{Cut elimination}

\begin{theorem}[Cut elimination]\label{thm:cut-elim}
Every derivation in the judgment 8-tuple system can be transformed into
a cut-free derivation.  The transformation preserves the coordinate,
proposition, and carrier, but may alter the evidence bundle and provenance.
Trust is non-increasing: $\trust' \tleq \trust$.
\end{theorem}

\begin{proof}[Proof sketch]
This follows from \Cref{thm:cut} by induction on the number of cut
applications in the derivation.  Each cut elimination step replaces
one cut with zero or more cuts of strictly smaller complexity, and the
double induction terminates by the well-foundedness argument given in
\S\ref{ssec:cut}.  The trust bound $\trust' \tleq \trust$ follows from
the conservative-join property of $\tjoin$: merging two evidence bundles
cannot exceed the minimum trust.
\end{proof}

\subsection{Trust monotonicity}

\begin{theorem}[Trust monotonicity under reduction]\label{thm:trust-mono}
If $\J \leadsto^* \J'$ (zero or more transitions), then the trust of $\J'$
is bounded below by the trust of the evidence accumulated along the path:
\[
  \trust(\J') \;\succeq\;
    \min\bigl(\trust(\J),\,
    \min_{t \in \mathrm{path}} \trust(\mathrm{new\_evidence}(t))\bigr).
\]
In particular, no transition can silently raise trust above the level
warranted by the evidence.
\end{theorem}

\begin{proof}
Each transition either (i)~leaves trust unchanged (lateral, fixpoint),
(ii)~raises trust by adding higher-level evidence (forward with solver
discharge), or (iii)~adds an obstruction and drops trust to $\Tcontra$
(interrupt).  In all cases, the trust is bounded by the minimum of the
incoming trust and the trust of any new evidence.  The bound follows by
induction on the length of the transition path.
\end{proof}

\begin{corollary}[No trust laundering]
It is impossible to construct a chain of transitions that starts at trust
level $\Tunver$, passes through $\Tcopilot$, and arrives at $\Tproof$
without introducing $\Tproof$-level evidence at some step.
\end{corollary}

\subsection{Information preservation}

\begin{theorem}[Information preservation]\label{thm:info-pres}
Let $\J = \Jud{\coord}{\prop}{\carrier}{\evid}{\oblig}{\obstr}{\trust}{\prov}$
be a well-formed judgment.  Then restriction, transport, and composition
each preserve all eight fields: the output judgment $\J'$ satisfies
$\J'.\mathsf{field} \neq \bot$ for every field $\mathsf{field} \in
\{\coord, \prop, \carrier, \evid, \oblig, \obstr, \trust, \prov\}$, where
$\bot$ denotes the trivial (empty/zero) value.  More precisely:
\begin{enumerate}[label=(\alph*)]
  \item \textbf{Restriction.}  $\J|_V$ retains a valid coordinate ($V$),
        a restricted proposition, a restricted carrier, a filtered evidence
        bundle with $|\evid|_V| \geq 1$ (at least the items relevant to $V$),
        the relevant obligations and obstructions, the original trust, and
        an extended provenance.
  \item \textbf{Transport.}  $f_*(\J)$ retains all eight fields by
        covariant push-forward; the coordinate changes to $V = f(\coord)$
        but all metadata is carried along.
  \item \textbf{Composition.}  $\J_1 \circ \J_2$ retains a coordinate
        (the meet $\coord_1 \sqcap \coord_2$), a conjoined proposition,
        a met carrier, a merged evidence bundle with
        $|\evid_1 \cup \evid_2| \geq \max(|\evid_1|, |\evid_2|)$,
        the union of obligations and obstructions, the conservative trust
        $\trust_1 \tjoin \trust_2$, and a merged provenance.
\end{enumerate}
\end{theorem}

\begin{proof}
We verify each operation preserves each field.

\emph{Restriction.}  By Definition~3.1 (\S\ref{ssec:restriction}),
$\J|_V = \Jud{V}{\prop|_V}{\carrier|_V}{\evid|_V}{\oblig|_V}
{\obstr|_V}{\trust}{\prov \cup \{f\}}$.
The coordinate $V$ is non-trivial by hypothesis (it is a valid object in
$\Site$).  The restricted proposition $\prop|_V$ is obtained by narrowing
free variables to those bound at $V$; since $V$ is a sub-coordinate of
$\coord$, and $\prop$ is well-formed at $\coord$ by (WF2), the restriction
is a syntactically well-formed proposition at $V$.
The filtered evidence $\evid|_V$ retains at least the evidence items whose
support includes $V$; since the original evidence has at least one item
(by well-formedness), and each item's support includes its coordinate's
ancestry chain, at least one item survives.
Trust is unchanged ($\trust' = \trust$).
Provenance is strictly extended ($|\prov'| > |\prov|$).

\emph{Transport.}  By Definition~3.3 (\S\ref{ssec:transport}),
$f_*(\J) = \Jud{V}{f_*(\prop)}{f_*(\carrier)}{f_*(\evid)}{f_*(\oblig)}
{f_*(\obstr)}{\trust}{\prov \cup \{f\}}$.
Each component is mapped covariantly along $f$; since $f$ is a morphism
in $\Site$ (hence non-degenerate), each push-forward produces a non-trivial
result.  The key point is that transport does not \emph{filter} evidence
(unlike restriction): all evidence items are carried to the new coordinate,
so $|f_*(\evid)| = |\evid|$.  Trust and provenance are preserved as in
restriction.

\emph{Composition.}  By Definition~3.5 (\S\ref{ssec:composition}),
$\J_1 \circ \J_2 = \Jud{\coord_1 \sqcap \coord_2}{\prop_1 \wedge \prop_2}
{\carrier_1 \sqcap \carrier_2}{\evid_1 \cup \evid_2}{\oblig_1 \cup \oblig_2}
{\obstr_1 \cup \obstr_2}{\trust_1 \tjoin \trust_2}{\prov_1 \otimes \prov_2}$.
The coordinate meet $\coord_1 \sqcap \coord_2$ exists because the judgments
share a boundary (precondition of composition).
The conjoined proposition is non-trivial whenever the inputs are.
The evidence union satisfies
$|\evid_1 \cup \evid_2| \geq \max(|\evid_1|, |\evid_2|) \geq 1$.
The trust $\trust_1 \tjoin \trust_2 = \min(\trust_1, \trust_2)$ is
non-trivial whenever both inputs have trust above $\Tcontra$.
Provenance is non-trivial because $\prov_1 \otimes \prov_2$ merges
two non-empty DAGs.

In all three cases, no field is reduced to its trivial value, so
information is preserved.
\end{proof}


%% ═══════════════════════════════════════════════════════════════════════
\section{Related Work and Conclusion}\label{sec:related}
%% ═══════════════════════════════════════════════════════════════════════

\subsection{Related work}

\paragraph{Dependent type theories.}
Martin-L\"of type theory~\cite{martinlof84}, the Calculus of Inductive
Constructions~\cite{coquand88}, and Homotopy Type Theory~\cite{hott}
provide the foundational notion of judgment that we extend.  Our 8-tuple
is a conservative extension: setting $\evid$, $\oblig$, $\obstr$,
$\trust$, and $\prov$ to their trivial values recovers the classical
3-component judgment (\S\ref{sec:recovering}).

\paragraph{Effect systems and refinement types.}
\Fstar{}~\cite{fstar} and Liquid Haskell~\cite{liquidhaskell} enrich
types with specifications, partially addressing the evidence gap.
Our evidence bundles and trust algebra go further by making the
\emph{source} and \emph{quality} of evidence first-class.

\paragraph{Program logics.}
Hoare logic~\cite{hoare69}, separation logic~\cite{reynolds02,iris},
and Iris~\cite{iris} provide powerful reasoning frameworks.
We show in \S\ref{sec:recovering} that Hoare triples embed as degenerate
8-tuples.  The key addition is the trust dimension, which allows mixing
proofs from different verification engines at different confidence levels.

\paragraph{Proof assistants.}
\LEAN{}~\cite{lean4}, \Coq{}~\cite{coq}, Isabelle/HOL~\cite{isabelle},
and \Dafny{}~\cite{dafny} implement various judgment forms.
None tracks trust gradation, evidence provenance, or first-class
obstructions.

\paragraph{Trusted computing bases.}
The notion of trust levels is related to the TCB stratification in
seL4~\cite{sel4} and CertiKOS~\cite{certikos}.
Our trust algebra provides a formal algebraic framework for the same
intuition.

\paragraph{Provenance in databases.}
Data provenance~\cite{buneman01,green07} tracks the lineage of query
results.  Our provenance component applies the same idea to proof
artifacts.

\subsection{Conclusion}

We have introduced the judgment 8-tuple
$\J = (\coord, \prop, \carrier, \evid, \oblig, \obstr, \trust, \prov)$
as the central algebraic object of Judgment Geometry.  The 8-tuple extends
the classical type-theoretic judgment with five new dimensions---evidence,
obligations, obstructions, trust, and provenance---while remaining a
conservative extension of existing systems.

The five algebraic operations (restriction, transport, composition,
comparison, join) give the space of judgments a rich categorical structure,
and the standard structural rules (weakening, contraction, exchange, cut)
are admissible.  The metatheory guarantees subject reduction, cut
elimination, and trust monotonicity.

The \JuGeo{} implementation demonstrates that the 8-tuple is not merely a
theoretical construct: all 7 trust algebra axioms are verified on 100
programs (700 axiom checks, 100\% pass rate), with individual algebraic
operations completing in 0.15--19.5\,\textmu{}s.  By retaining all eight components
through every operation---a $2.0\times$ information advantage over the
best competing system---the 8-tuple enables new
verification workflows---progressive verification, trust-aware composition,
and obstruction-driven repair---that were previously impossible.

Future work will develop the sheaf-theoretic aspects of Judgment Geometry
(descent, gluing, cohomological obstruction theory) and scale the
evaluation to industrial codebases.

%% ═══════════════════════════════════════════════════════════════════════
%% Bibliography
%% ═══════════════════════════════════════════════════════════════════════

\begin{thebibliography}{99}

\bibitem{martinlof84}
P.~Martin-L\"of.
\newblock \emph{Intuitionistic Type Theory}.
\newblock Bibliopolis, Naples, 1984.

\bibitem{coquand88}
T.~Coquand and C.~Huet.
\newblock The {Calculus of Constructions}.
\newblock \emph{Information and Computation}, 76(2--3):95--120, 1988.

\bibitem{gentzen35}
G.~Gentzen.
\newblock Untersuchungen \"uber das logische {S}chlie{\ss}en.
\newblock \emph{Mathematische Zeitschrift}, 39:176--210, 405--431, 1935.

\bibitem{hott}
The Univalent Foundations Program.
\newblock \emph{Homotopy Type Theory: Univalent Foundations of Mathematics}.
\newblock Institute for Advanced Study, 2013.

\bibitem{lean4}
L.~de~Moura and S.~Ullrich.
\newblock The {Lean~4} theorem prover and programming language.
\newblock In \emph{CADE-28}, LNCS 12699, pp.~625--635. Springer, 2021.

\bibitem{coq}
The Coq Development Team.
\newblock \emph{The {Coq} Reference Manual}, version 8.18, 2023.

\bibitem{fstar}
N.~Swamy, C.~Hri\c{t}cu, C.~Keller, A.~Rastogi, et~al.
\newblock Dependent types and multi-monadic effects in {F*}.
\newblock In \emph{POPL 2016}, pp.~256--270. ACM, 2016.

\bibitem{dafny}
K.~R.~M. Leino.
\newblock Dafny: An automatic program verifier for functional correctness.
\newblock In \emph{LPAR-16}, LNCS 6355, pp.~348--370. Springer, 2010.

\bibitem{hoare69}
C.~A.~R. Hoare.
\newblock An axiomatic basis for computer programming.
\newblock \emph{Communications of the ACM}, 12(10):576--580, 1969.

\bibitem{reynolds02}
J.~C. Reynolds.
\newblock Separation logic: A logic for shared mutable data structures.
\newblock In \emph{LICS 2002}, pp.~55--74. IEEE, 2002.

\bibitem{iris}
R.~Jung, R.~Krebbers, J.-H.~Jourdan, A.~Bizjak, L.~Birkedal, and D.~Dreyer.
\newblock Iris from the ground up: A modular foundation for higher-order
  concurrent separation logic.
\newblock \emph{Journal of Functional Programming}, 28:e20, 2018.

\bibitem{isabelle}
T.~Nipkow, L.~C. Paulson, and M.~Wenzel.
\newblock \emph{Isabelle/HOL --- A Proof Assistant for Higher-Order Logic}.
\newblock LNCS 2283. Springer, 2002.

\bibitem{liquidhaskell}
N.~Vazou, E.~L. Seidel, R.~Jhala, D.~Vytiniotis, and S.~Peyton~Jones.
\newblock Refinement types for {Haskell}.
\newblock In \emph{ICFP 2014}, pp.~269--282. ACM, 2014.

\bibitem{sel4}
G.~Klein, K.~Elphinstone, G.~Heiser, et~al.
\newblock {seL4}: Formal verification of an {OS} kernel.
\newblock In \emph{SOSP 2009}, pp.~207--220. ACM, 2009.

\bibitem{certikos}
R.~Gu, Z.~Shao, H.~Chen, X.~Wu, J.~Kim, V.~Sj\"oberg, and D.~Costanzo.
\newblock {CertiKOS}: An extensible architecture for building certified
  concurrent {OS} kernels.
\newblock In \emph{OSDI 2016}, pp.~653--669. USENIX, 2016.

\bibitem{buneman01}
P.~Buneman, S.~Khanna, and W.-C. Tan.
\newblock Why and where: A characterization of data provenance.
\newblock In \emph{ICDT 2001}, LNCS 1973, pp.~316--330. Springer, 2001.

\bibitem{green07}
T.~J. Green, G.~Karvounarakis, and V.~Tannen.
\newblock Provenance semirings.
\newblock In \emph{PODS 2007}, pp.~31--40. ACM, 2007.

\bibitem{Girard1989}
J.-Y.~Girard.
\newblock Geometry of interaction~I: Interpretation of {System~F}.
\newblock In \emph{Logic Colloquium '88}, pages 221--260. North-Holland, 1989.

\end{thebibliography}

\end{document}
